% ==================================================================================================================================

%%%%%%%%%%%%%%%%%%%%%%%%%%%
\chapter{Trabalhos Relacionados}
\label{trabRelac}
%%%%%%%%%%%%%%%%%%%%%%%%%%%

Este capítulo apresenta os trabalhos relacionados às duas vertentes que fundamentam esta pesquisa. A primeira, apresentada na Seção~\ref{sec:fs-ca}, compreende estudos que investigam a utilização de métodos de seleção de atributos em classificadores associativos. A segunda, apresentada na Seção~\ref{ensambleFS}, aborda estratégias de \textit{Ensemble Feature Selection} (EFS), com ênfase na diversidade entre seletores, nos mecanismos de agregação, na estabilidade da seleção e na definição dos subconjuntos finais de atributos. Por fim, a Seção~\ref{sec:consideracoes-relacionados} sintetiza as evidências identificadas e as lacunas que motivam esta pesquisa.

% =====================================================================
% =====================================================================
% =====================================================================
% =====================================================================
\section{Seleção de Atributos em Classificadores Associativos}
\label{sec:fs-ca}
% =====================================================================

A seleção de atributos assume uma função
particularmente relevante em classificação associativa, uma vez que a
dimensionalidade do conjunto de dados influencia não apenas a etapa de
classificação, mas também o espaço de busca utilizado para a descoberta de
\textit{itemsets} e Regras de Associação de Classe (\textit{Class Association
Rules} -- CARs). A presença de atributos irrelevantes ou redundantes pode
aumentar o número de combinações candidatas, favorecer a geração de regras
redundantes e elevar o custo computacional da mineração. Consequentemente,
reduzir o espaço de atributos pode afetar simultaneamente a complexidade do
processo de mineração, a quantidade de regras produzidas, a interpretabilidade
do modelo e seu desempenho preditivo.

A literatura apresenta diferentes estratégias para incorporar FS a
classificadores associativos. Em algumas propostas, a seleção é realizada como
uma etapa independente de pré-processamento, de forma que somente o subconjunto
selecionado é submetido posteriormente à mineração de regras
\citeonline{piri2019exploring,alhawari2021classification,sowan2023hybrid,
alagukumar2025impacts}. Em outras, a relevância dos atributos participa
diretamente da geração ou ponderação das CARs
\citeonline{nandhini2022performance,wu2024novel}. Há ainda propostas nas quais
a própria estrutura das regras é utilizada como fonte de informação para
selecionar características, estabelecendo uma integração mais estreita entre
FS e classificação associativa \citeonline{rad2024framework}.

Diante dessa diversidade, esta seção analisa os trabalhos que relacionam
explicitamente FS e classificação associativa, buscando identificar como a
seleção é inserida no \textit{pipeline}, quais critérios são utilizados para
avaliar atributos, quais classificadores associativos são empregados e quais
efeitos são observados sobre desempenho, dimensionalidade e estrutura das
regras.

% ---------------------------------------------------------------------
\subsection{Procedimento da revisão}
\label{subsec:procedimento-rsl-ca}
% ---------------------------------------------------------------------

A identificação dos estudos foi conduzida por meio de uma revisão da literatura
orientada especificamente à relação entre FS e classificação associativa. A
estratégia de busca combinou termos associados à classificação baseada em regras,
incluindo \textit{associative classification}, \textit{associative classifier},
\textit{class association rules} e \textit{predictive association rules}, com
termos relacionados à redução de atributos, como \textit{feature selection},
\textit{attribute selection}, \textit{feature ranking} e \textit{item
selection}. As consultas foram adaptadas às particularidades das bases
bibliográficas utilizadas no levantamento.

A análise originalmente foi orientada por cinco questões de pesquisa, voltadas
a: (i) identificar o \textit{pipeline} adotado e as etapas modificadas ou
preservadas; (ii) verificar se FS ocorre externamente ou é incorporada ao
classificador; (iii) caracterizar os métodos de seleção utilizados; (iv)
identificar estratégias adicionais de redução de complexidade ou melhoria do
modelo; e (v) examinar a avaliação experimental apresentada pelos estudos.
Nesta seção, essas questões são consolidadas em três eixos analíticos:
\textit{forma de integração}, \textit{abordagens de seleção} e
\textit{evidências experimentais}, evitando a repetição da descrição
metodológica da revisão.

Após a aplicação dos critérios de seleção e a atualização do levantamento, o conjunto de trabalhos analisado nesta seção compreende 12 estudos que apresentam relação
explícita entre seleção ou avaliação de atributos e classificadores
associativos. Os trabalhos são discutidos a seguir em ordem alfabética pelo
sobrenome do primeiro autor.

% ---------------------------------------------------------------------
\subsection{Trabalhos analisados}
\label{subsec:trabalhos-ca}
% ---------------------------------------------------------------------

\noindent\textbf{\citeonline{abubacker2017integrated}.}
O trabalho propõe uma arquitetura híbrida para classificação de mamografias na
qual a redução dimensional antecede a mineração associativa. Inicialmente,
\citeonline{abubacker2017integrated} empregam um mecanismo de seleção
multivariada baseado em \textit{Correlation-based Feature Selection} (CFS),
associado a uma busca genética, para obter um subconjunto de características
relevantes. Os atributos selecionados são então fornecidos a uma versão
modificada do GARM (\textit{Genetic Association Rule Mining}), responsável pela
geração das regras. Posteriormente, essas regras são utilizadas como entrada
para classificadores neurais e neuro-fuzzy, denominados ACNN e ACFNN. A
estrutura proposta evidencia uma integração em estágios, na qual FS restringe
previamente o espaço submetido à mineração de regras. Experimentalmente, o ACFNN
alcançou acurácia média superior à versão neural, indicando benefício da
combinação entre seleção, mineração associativa e sistema neuro-fuzzy.

\medskip

\noindent\textbf{\citeonline{alagukumar2025impacts}.}
A relação entre FS e discretização é investigada explicitamente por
\citeonline{alagukumar2025impacts} em dados de expressão gênica. Os autores
avaliam três filtros baseados em informação mútua --- JIM
(\textit{Joint Impurity Filter}), JMI (\textit{Joint Mutual Information}) e JMI3
(\textit{Third-order Joint Mutual Information}) --- combinados com cinco
métodos supervisionados de discretização: MDLP, CAIM, CACC, AMEVA e
Chi-Merge. O subconjunto resultante é discretizado, transformado para o formato
transacional e posteriormente submetido aos classificadores CBA e CMAR. Dessa
forma, o estudo não trata FS como uma decisão isolada, mas investiga seu efeito
em conjunto com a granularidade criada pela discretização.

Além do desempenho preditivo, os autores analisam o
número de intervalos, \textit{itemsets} e o tempo de execução. Nos conjuntos de
expressão gênica considerados, o CMAR atingiu 100\% de acurácia nas combinações
avaliadas, enquanto o comportamento do CBA apresentou maior dependência da
combinação entre filtro e discretizador. Esse resultado é relevante para
classificação associativa porque evidencia que FS e discretização afetam
diretamente a representação transacional sobre a qual as CARs são geradas.

\medskip

\noindent\textbf{\citeonline{alazaidah2021feature}.}
Em um estudo comparativo especificamente direcionado à classificação
associativa, \citeonline{alazaidah2021feature} avaliam nove métodos de seleção:
ABB-IEP, ABB-LIU, Focus, LVF, Relief, SBS-LIU, SBS-IEP, SFS-IEP e SFS-LIU.
Os subconjuntos produzidos são avaliados com cinco classificadores
associativos: CBA, CBA2, CFAR, CMAR e CPAR. A comparação é realizada sobre
conjuntos de dados com diferentes tipos de atributos, permitindo analisar não
apenas qual seletor produz maior acurácia, mas também sua interação com a
natureza dos dados e com o classificador utilizado.

Os resultados de \citeonline{alazaidah2021feature} mostram que não há uma
configuração única dominante em todos os cenários: ABB-IEP e ABB-LIU se
destacam nos conjuntos compostos por atributos nominais, enquanto Relief
apresenta melhor comportamento em conjuntos com atributos reais e inteiros.
Esse resultado reforça a dependência entre o mecanismo de seleção e as
características do problema, aspecto importante quando se considera a adoção de
um único seletor como etapa fixa de um \textit{pipeline} associativo.

\medskip

\noindent\textbf{\citeonline{alhawari2021classification}.}
Investigam classificação de avaliações
de aplicativos em categorias relacionadas à manutenção de software. Após etapas
de processamento de linguagem natural e transformação dos textos em uma
representação estruturada, os autores aplicam \textit{Information Gain} (IG) e
$\chi^2$ para selecionar os atributos textuais mais relevantes. Os subconjuntos
obtidos são utilizados com classificadores tradicionais e associativos,
incluindo CBA, CBA2, CMAR e CPAR.

O trabalho também propõe o ACRM (\textit{Associative Classification for Reviews
Mining}), que utiliza FP-Growth para obtenção de padrões e combina confiança e
\textit{conviction} na seleção das regras. Na análise de FS,
\citeonline{alhawari2021classification} verificam melhoria de desempenho após a
remoção de atributos pouco informativos, com IG apresentando, em geral, melhores
resultados que $\chi^2$. Os autores observam, entretanto, que os classificadores
associativos são menos sensíveis à presença de atributos ineficientes do que
alguns classificadores tradicionais, em razão da utilização das regras mais
fortes no processo de predição.

\medskip

\noindent\textbf{\citeonline{ashi2021enhancing}.}
No domínio de predição de \textit{customer churn},
analisam o efeito da FS sobre CBA, CBA2, CPAR e
FARC-HD. Após o pré-processamento dos dados, é utilizado o SBS-IEP
(\textit{Sequential Backward Search using Inconsistent Examples Pairs}) para
reduzir o conjunto de atributos antes da indução dos classificadores
associativos. O estudo compara os modelos com e sem seleção por meio de
acurácia, taxa de erro, curva ROC, quantidade de regras e tempo de execução.

A contribuição do estudo de \citeonline{ashi2021enhancing} é particularmente
relevante por avaliar não apenas desempenho preditivo, mas também efeitos
estruturais da redução dimensional. Os resultados mostram que a aplicação de FS
altera a quantidade de regras produzidas e o custo computacional dos
classificadores, evidenciando que, em classificação associativa, a seleção de
atributos pode repercutir diretamente sobre a complexidade do modelo final.

\medskip

\noindent\textbf{\citeonline{cicek2021assessment}.}
Aplicam classificação associativa a dados de
expressão gênica para distinguir indivíduos positivos e negativos para COVID-19.
O conjunto original contém 15.979 genes de 234 indivíduos, caracterizando um
problema no qual a dimensionalidade é substancialmente superior ao número de
amostras. Para reduzir esse espaço, os autores utilizam LASSO como mecanismo de
seleção dos preditores e, posteriormente, constroem o modelo de classificação
associativa sobre o subconjunto selecionado.

A avaliação reportada por \citeonline{cicek2021assessment} inclui acurácia,
acurácia balanceada, sensibilidade, especificidade, valor preditivo positivo,
valor preditivo negativo e F1-\textit{score}. O modelo associativo alcança
92,7\% de acurácia e 90,5\% de F1-\textit{score}. Além da redução de um espaço
inicial de milhares de genes, o trabalho utiliza as regras resultantes para
relacionar combinações de genes às classes, preservando a finalidade
interpretável do classificador associativo.

\medskip

\noindent\textbf{\citeonline{ghareb2018enhanced}.}
No contexto de categorização de textos em língua árabe,
os autores propõem três métodos de filtro, denominados
CRFM, MCDM e OR2. Os métodos são comparados com doze estratégias tradicionais
de seleção, totalizando uma avaliação ampla de diferentes critérios de FS. Os
subconjuntos resultantes são avaliados tanto com Naive Bayes quanto com
classificação associativa, sobre três conjuntos de documentos.

Os resultados de \citeonline{ghareb2018enhanced} mostram que os filtros
propostos podem produzir desempenho competitivo ou superior em diferentes
configurações, porém sem identificar um único método dominante para todos os
conjuntos de dados, tamanhos de subconjunto e classificadores. Esse resultado é
importante para a presente discussão porque fornece evidência empírica de que
diferentes critérios de FS capturam propriedades distintas dos dados, limitando
a adoção de um seletor individual como solução universal.

\medskip

\noindent\textbf{\citeonline{nandhini2022performance}.}
Diferentemente dos trabalhos que realizam FS como pré-processamento,
\citeonline{nandhini2022performance} investigam a seleção de atributos dentro do
próprio processo de geração de regras do CPAR. O algoritmo original utiliza
FOIL's Information Gain para selecionar, a cada iteração, o atributo que será
incorporado ao antecedente das CARs. Os autores substituem esse critério por
\textit{Gain Ratio} (GR) e \textit{Mutual Information Gain} (IG), originando as
variantes CPAR-GR e CPAR-IG.

Além do critério de seleção, \citeonline{nandhini2022performance} avaliam
\textit{Laplace Accuracy} e \textit{Likelihood Ratio} na avaliação e escolha das
melhores regras. Os experimentos são conduzidos em diferentes conjuntos de
dados médicos, considerando acurácia e AUC e complementados por testes
estatísticos. Os resultados indicam comportamento particularmente favorável do
CPAR-GR, mostrando que a escolha da medida utilizada durante a construção das
regras altera diretamente o desempenho do classificador. Assim, o estudo
constitui um exemplo de FS \textit{interna}, na qual seleção de atributos e
geração de CARs não são etapas independentes.

\medskip

\noindent\textbf{\citeonline{piri2019exploring}.}
Utilizam classificação associativa para determinar
o estado fetal a partir de dados de cardiotocografia. Os atributos contínuos são
inicialmente transformados por discretização baseada em entropia e, em seguida,
o \textit{ExtraTreesClassifier} é empregado para calcular a importância das
variáveis. A partir desse ranking são selecionados dez atributos, que compõem a
versão reduzida do conjunto de dados utilizada pelo classificador CBA.

O estudo compara o CBA antes e depois de FS e reporta aumento de acurácia de
83\% para 84\%. Além disso, \citeonline{piri2019exploring} observam redução no
número de regras e no tempo necessário para construção do classificador após a
remoção dos atributos menos relevantes. Embora o ganho preditivo seja limitado,
o resultado evidencia os efeitos simultâneos da seleção sobre desempenho e
compacidade do modelo associativo.

\medskip

\noindent\textbf{\citeonline{rad2024framework}.}
Propõem um \textit{framework} para detecção de
moedas falsificadas baseado em classificação associativa fuzzy. A proposta
introduz o PFA (\textit{Pruned Fuzzy Associative Classifier}) e o mecanismo
Engine\_FAFS, no qual as próprias regras de associação fuzzy são utilizadas para
identificar características relevantes. Nesse caso, FS não é simplesmente
executada como um filtro independente antes da mineração; a evidência obtida das
regras participa da determinação do subconjunto empregado pelo sistema.

O \textit{framework} de \citeonline{rad2024framework} também combina seleção de
características, poda de regras e otimização por \textit{Particle Swarm
Optimization} (PSO). Entre os parâmetros otimizados encontra-se o limiar
$\delta$, relacionado ao tamanho do subconjunto selecionado. A proposta,
portanto, representa uma forma de integração baseada nas próprias regras,
distinta tanto do pré-processamento externo quanto da seleção dinâmica durante a
construção dos antecedentes.

\medskip

\noindent\textbf{\citeonline{sowan2023hybrid}.}
Um dos trabalhos mais diretamente relacionados à combinação de múltiplos
seletores é apresentado por \citeonline{sowan2023hybrid}, que propõem duas
estratégias de FS para classificação associativa de câncer de mama. A primeira,
denominada WARF (\textit{Weighted Average of Relative Frequency}), constitui um
\textit{ensemble} de filtros que agrega CFS, \textit{Information Gain},
\textit{Gain Ratio}, $\chi^2$ e ReliefF. A seleção combina os resultados desses
métodos por meio da frequência relativa de escolha dos atributos, produzindo um
subconjunto que é posteriormente utilizado pelo PCBA.

A segunda estratégia de \citeonline{sowan2023hybrid}, denominada Hybrid
PSO-WARF, acrescenta ao resultado do ensemble uma etapa de busca baseada em
\textit{Particle Swarm Optimization}, formando uma abordagem híbrida entre
filtro e \textit{wrapper}. Os autores avaliam WARF-PCBA e Hybrid
PSO-WARF-PCBA em dois conjuntos de câncer de mama, comparando-os com CBA e com
configurações individuais baseadas em ReliefF, CFS, IG, GR e $\chi^2$, além de
outros classificadores. No conjunto de recorrência de câncer de mama, por
exemplo, a utilização de WARF elevou a acurácia em relação ao conjunto completo
de atributos, e a variante híbrida atingiu 75,9\%.

Entre os estudos analisados nesta seção, \citeonline{sowan2023hybrid}
constituem a evidência mais explícita de aplicação de
\textit{Ensemble Feature Selection} diretamente associada a um classificador
associativo. Por essa razão, o trabalho é especialmente relevante para a
discussão desenvolvida na Seção~\ref{ensambleFS}, embora sua estratégia de
agregação e sua etapa híbrida sejam distintas da abordagem investigada nesta
dissertação.

\medskip

\noindent\textbf{\citeonline{wu2024novel}.}
\citeonline{wu2024novel} propõem o MCWCAR (\textit{Mutual Information and
Correlation Coefficient Weighted Class Association Rule Mining}) para predição
de defeitos de software. A principal diferença em relação aos métodos de
pré-processamento é que a relevância das características é incorporada à própria
mineração dos \textit{itemsets}. A informação mútua é utilizada na ponderação
das características, enquanto medidas de correlação contribuem para distinguir
relações associadas às classes de defeito e não defeito.

O MCWCAR integra, portanto, seleção/ponderação de atributos e mineração de CARs
em um único mecanismo. \citeonline{wu2024novel} avaliam a proposta em 27
conjuntos de dados do repositório PROMISE e utilizam métricas como
\textit{Balance}, MCC, G-Mean e F-\textit{measure}. As comparações com CWCAR,
CBA3 e ACAR são complementadas por Scott-Knott ESD e Wilcoxon
\textit{signed-rank}, sendo observadas diferenças estatisticamente
significativas em favor do MCWCAR para diferentes métricas. O estudo demonstra
que incorporar critérios tradicionalmente relacionados a FS à própria mineração
pode reduzir o espaço de busca e simultaneamente alterar a qualidade das regras
geradas.

% ---------------------------------------------------------------------
\subsection{Síntese comparativa}
\label{subsec:sintese-ca}
% ---------------------------------------------------------------------

A análise dos trabalhos mostra que FS é utilizada em diferentes níveis de
acoplamento com a classificação associativa. A Tabela~\ref{tab:sintese-ca}
sintetiza os 12 estudos segundo a forma de integração, o mecanismo de seleção
utilizado, o classificador associativo empregado e a principal característica
metodológica. As referências são apresentadas em ordem alfabética.

\begingroup
\scriptsize
\setlength{\tabcolsep}{3pt}
\renewcommand{\arraystretch}{1.15}

\begin{longtable}{@{}p{2.65cm}p{1.65cm}p{3.65cm}p{2.6cm}p{4.95cm}@{}}

\caption{Síntese dos trabalhos que integram seleção de atributos e classificação associativa.}
\label{tab:sintese-ca}\\

\toprule
\textbf{Referência} &
\textbf{Integração} &
\textbf{Método de FS} &
\textbf{CA} &
\textbf{Característica principal} \\
\midrule
\endfirsthead

\multicolumn{5}{c}{\tablename\ \thetable{} -- continuação}\\
\toprule
\textbf{Referência} &
\textbf{Integração} &
\textbf{Método de FS} &
\textbf{CA} &
\textbf{Característica principal} \\
\midrule
\endhead

\midrule
\multicolumn{5}{r}{\textit{continua na próxima página}}\\
\endfoot

\bottomrule
\endlastfoot

\citeonline{abubacker2017integrated} &
Externa &
CFS com busca genética &
GARM + ACNN/ACFNN &
Seleciona características de mamografia antes da mineração de regras; as regras
geradas posteriormente alimentam classificadores neural e neuro-fuzzy. \\

\citeonline{alagukumar2025impacts} &
Externa &
JIM, JMI e JMI3 &
CBA e CMAR &
Analisa conjuntamente FS baseada em informação mútua e cinco métodos de
discretização, relacionando seleção ao número de intervalos,
\textit{itemsets}, tempo e desempenho. \\

\citeonline{alazaidah2021feature} &
Externa &
ABB-IEP, ABB-LIU, Focus, LVF, Relief, SBS e SFS &
CBA, CBA2, CFAR, CMAR e CPAR &
Compara nove seletores com cinco CAs e mostra dependência entre método de FS,
tipo de atributo e classificador. \\

\citeonline{alhawari2021classification} &
Externa &
IG e $\chi^2$ &
CBA, CBA2, CMAR, CPAR e ACRM &
Reduz a representação textual de avaliações de aplicativos e investiga o efeito
da seleção sobre classificadores tradicionais e associativos. \\

\citeonline{ashi2021enhancing} &
Externa &
SBS-IEP &
CBA, CBA2, CPAR e FARC-HD &
Avalia modelos com e sem FS considerando desempenho, número de regras e tempo
de execução. \\

\citeonline{cicek2021assessment} &
Externa &
LASSO &
Classificação associativa/CBA &
Reduz 15.979 genes antes da indução do modelo associativo para classificação de
COVID-19 e interpretação das regras gênicas. \\

\citeonline{ghareb2018enhanced} &
Externa &
CRFM, MCDM e OR2 &
AC &
Propõe três filtros e os compara com doze métodos tradicionais em categorização
de textos árabes, sem identificar um seletor universalmente dominante. \\

\citeonline{nandhini2022performance} &
Interna &
FOIL-IG, GR e IG &
CPAR e variantes &
Seleciona dinamicamente atributos durante a formação dos antecedentes das CARs,
originando CPAR-GR e CPAR-IG. \\

\citeonline{piri2019exploring} &
Externa &
ExtraTreesClassifier &
CBA &
Seleciona dez atributos após discretização por entropia; compara CBA com e sem
FS e observa redução de regras e tempo. \\

\citeonline{rad2024framework} &
Integrada às regras &
Engine\_FAFS + PSO &
PFA &
Utiliza as próprias regras fuzzy para identificar características relevantes e
combina FS com otimização e poda de regras. \\

\citeonline{sowan2023hybrid} &
Externa/híbrida &
WARF: CFS + IG + GR + $\chi^2$ + ReliefF; PSO-WARF &
PCBA &
Emprega explicitamente ensemble heterogêneo de filtros e uma variante híbrida
com PSO antes da classificação associativa. \\

\citeonline{wu2024novel} &
Interna &
Informação mútua e coeficiente de correlação &
MCWCAR &
Incorpora relevância dos atributos à mineração por meio de pesos utilizados na
descoberta e avaliação dos \textit{itemsets}. \\

\end{longtable}
\endgroup

% ---------------------------------------------------------------------
\subsection{Análise dos estudos}
\label{subsec:analise-ca}
% ---------------------------------------------------------------------

\subsubsection{Forma de integração entre FS e classificação associativa}

A primeira distinção identificada refere-se ao momento em que a seleção de
atributos interfere no processo associativo. Em nove dos 12 estudos, a FS é
realizada antes da indução do classificador, caracterizando uma integração
externa. Essa configuração é observada em
\citeonline{abubacker2017integrated},
\citeonline{alagukumar2025impacts},
\citeonline{alazaidah2021feature},
\citeonline{alhawari2021classification},
\citeonline{ashi2021enhancing},
\citeonline{cicek2021assessment},
\citeonline{ghareb2018enhanced},
\citeonline{piri2019exploring} e
\citeonline{sowan2023hybrid}. Apesar das diferenças entre os seletores, esses
trabalhos compartilham uma mesma separação conceitual: primeiro é determinado um
subconjunto de atributos e, posteriormente, o classificador associativo é
construído sobre o espaço reduzido.

Essa separação apresenta como vantagem a modularidade, pois permite modificar o
seletor sem necessariamente alterar o algoritmo responsável pela mineração das
CARs. Tal característica é particularmente evidente nos experimentos de
\citeonline{alazaidah2021feature}, nos quais diferentes seletores são combinados
com diferentes CAs, e em \citeonline{sowan2023hybrid}, que compara seletores
individuais, o ensemble WARF e sua variante híbrida mantendo a etapa associativa
subsequente.

Em contraste, \citeonline{nandhini2022performance} e
\citeonline{wu2024novel} integram a informação de relevância diretamente ao
processo de indução. No primeiro caso, o atributo incorporado ao antecedente é
escolhido dinamicamente durante a geração de cada regra; no segundo, pesos
baseados em informação mútua e correlação modificam a mineração dos
\textit{itemsets}. Assim, a FS deixa de ser uma transformação independente do
conjunto de entrada e passa a interferir diretamente no espaço de regras
explorado pelo algoritmo.

A proposta de \citeonline{rad2024framework} ocupa uma posição distinta. O
Engine\_FAFS utiliza informação proveniente das próprias regras fuzzy para
decidir quais características devem permanecer no modelo. Portanto, o trabalho
não se enquadra adequadamente em uma divisão estritamente binária entre seleção
``antes'' ou ``durante'' a mineração. Trata-se de uma integração orientada pelas
regras, combinada posteriormente com poda e otimização por PSO.

\subsubsection{Abordagens de seleção e relação com a geração de regras}

Os 12 estudos mostram elevada diversidade nos critérios empregados para
representar relevância. Medidas baseadas em teoria da informação aparecem de
forma recorrente. \citeonline{nandhini2022performance} exploram FOIL-IG, GR e
IG dentro do CPAR; \citeonline{alhawari2021classification} utilizam IG em
conjunto com $\chi^2$; \citeonline{alagukumar2025impacts} analisam JIM, JMI e
JMI3; e \citeonline{wu2024novel} incorporam informação mútua à ponderação dos
atributos. Essa recorrência indica que relações entre atributo e classe
constituem uma das principais fontes de evidência utilizadas na literatura para
controlar o espaço de busca associativo.

Entretanto, os trabalhos também mostram que essa evidência não é suficiente para
definir um único critério preferencial. \citeonline{alazaidah2021feature}
identificam comportamentos distintos conforme o tipo dos atributos,
\citeonline{ghareb2018enhanced} observam variação entre seletores, conjuntos de
dados e classificadores, e \citeonline{nandhini2022performance} mostram que a
substituição da medida utilizada pelo CPAR altera significativamente o
desempenho das regras geradas. Em conjunto, esses resultados indicam que
diferentes métodos de FS capturam propriedades distintas da relação entre os
atributos e a variável-alvo.

Também são encontrados métodos dependentes de modelos ou mecanismos de busca.
\citeonline{cicek2021assessment} utilizam LASSO em um cenário de elevada
dimensionalidade; \citeonline{piri2019exploring} empregam a importância
calculada pelo ExtraTreesClassifier; \citeonline{abubacker2017integrated}
combinam CFS com busca genética; e \citeonline{ashi2021enhancing} utilizam
busca sequencial para eliminação de atributos. Essas estratégias reforçam a
heterogeneidade metodológica existente na literatura.

O caso de \citeonline{sowan2023hybrid} diferencia-se dos demais porque a
heterogeneidade dos critérios é explorada deliberadamente. O WARF não escolhe
entre CFS, IG, GR, $\chi^2$ e ReliefF; em vez disso, agrega a evidência dos
cinco seletores para determinar os atributos mais recorrentes. A variante
Hybrid PSO-WARF acrescenta ainda uma busca orientada pelo desempenho. Dessa
forma, o trabalho transforma a discordância entre seletores em parte explícita
do mecanismo de seleção, característica central de uma abordagem de
\textit{Ensemble Feature Selection}.

\subsubsection{Efeitos da FS e avaliação experimental}

A avaliação dos estudos indica que os efeitos de FS não devem ser interpretados
somente em termos de acurácia. Em problemas associativos, a redução do espaço de
atributos pode modificar o número de \textit{itemsets} candidatos, a quantidade
de CARs, o tempo de mineração e a complexidade do classificador resultante.
Essa preocupação é explicitamente observada por
\citeonline{ashi2021enhancing},
\citeonline{piri2019exploring},
\citeonline{alagukumar2025impacts} e
\citeonline{sowan2023hybrid}, que analisam aspectos adicionais à capacidade
preditiva.

\citeonline{alagukumar2025impacts}, por exemplo, mostram que a escolha do
seletor deve ser analisada juntamente com a discretização, uma vez que os
intervalos produzidos alteram a representação transacional e o número de
\textit{itemsets}. De maneira complementar,
\citeonline{piri2019exploring} associam a remoção de atributos à redução da
quantidade de regras e do tempo de construção, enquanto
\citeonline{ashi2021enhancing} reportam alterações tanto no número de regras
quanto no custo de execução após FS.

Em termos de rigor experimental, os estudos também apresentam diferenças.
\citeonline{nandhini2022performance} utilizam testes estatísticos para comparar
as variantes do CPAR, \citeonline{alhawari2021classification} complementam suas
comparações com testes de significância e tamanho de efeito, e
\citeonline{wu2024novel} utilizam Scott-Knott ESD e Wilcoxon
\textit{signed-rank} na avaliação sobre 27 conjuntos de dados. Por outro lado,
parte dos estudos concentra a avaliação em um domínio específico e em um número
mais restrito de bases, como ocorre em
\citeonline{abubacker2017integrated},
\citeonline{cicek2021assessment},
\citeonline{piri2019exploring} e
\citeonline{rad2024framework}. Essas diferenças limitam comparações diretas
entre os valores absolutos reportados pelos trabalhos.

De forma geral, os resultados disponíveis sustentam a utilização de FS como
mecanismo de controle da dimensionalidade em classificação associativa, mas não
permitem estabelecer um seletor individual como superior de maneira
independente do conjunto de dados e do classificador. Essa conclusão é
particularmente apoiada pelas comparações de
\citeonline{alazaidah2021feature} e
\citeonline{ghareb2018enhanced}, nas quais o melhor comportamento varia entre
configurações, bem como pelos resultados de
\citeonline{nandhini2022performance}, que demonstram a influência do critério
utilizado dentro de um mesmo classificador.

Nesse contexto, a abordagem de \citeonline{sowan2023hybrid} constitui um ponto
de ligação importante com a literatura de \textit{Ensemble Feature Selection},
pois substitui a dependência de um único critério pela agregação de diferentes
seletores antes da classificação associativa. Entretanto, entre os estudos
analisados, estratégias desse tipo permanecem pouco exploradas quando comparadas
ao uso de seletores individuais. Essa constatação motiva a análise específica
de métodos de \textit{Ensemble Feature Selection} apresentada na
Seção~\ref{ensambleFS}, com o objetivo de investigar como diferentes evidências de
relevância podem ser combinadas de maneira sistemática antes de sua integração
com classificadores associativos.

% =====================================================================
% 3.2 -- ENSEMBLE FEATURE SELECTION
% =====================================================================
%%%%%%%%%%%%%%%%%%%%%%%%%%%
\section{\textit{Ensemble Feature Selection}}\label{ensambleFS}
%%%%%%%%%%%%%%%%%%%%%%%%%%%

O \textit{Ensemble Feature Selection} (EFS) combina resultados de diferentes métodos de seleção de atributos para reduzir a dependência de um único critério de avaliação e explorar a complementaridade entre os seletores \citeonline{bolon2019ensembles}. Sua construção envolve decisões relacionadas à composição do \textit{ensemble}, à diversidade entre os métodos, à representação e agregação dos resultados e à definição do subconjunto final de atributos. Essas decisões podem influenciar tanto a seleção obtida quanto sua estabilidade e capacidade de generalização, conforme discutido por \citeonline{akhy2024comprehensive}.

Esta seção apresenta e analisa 22 estudos primários sobre EFS, identificados por meio de buscas em bases bibliográficas (\textit{Scopus}, \textit{IEEE Xplore}, \textit{ACM Digital Library}, \textit{Web of Science} e \textit{Google Scholar}), utilizando combinações de termos relacionados a \textit{ensemble feature selection} e agregação de \textit{rankings}\footnote{A expressão de referência utilizada na Scopus foi \textit{TITLE((ensemble* W/2 ``feature selection'') OR (``rank aggregation'' W/2 ``feature selection''))}, com adaptações à sintaxe das demais bases consultadas.}. Trabalhos de revisão foram utilizados como apoio conceitual, mas não integram a comparação metodológica. A Tabela~\ref{tab:efs} sintetiza tais estudos, considerando a composição do \textit{ensemble}, a representação dos resultados, as estratégias de agregação, os critérios adicionais empregados, a definição do subconjunto final e as formas de avaliação. Com base nessa caracterização, as subseções seguintes discutem comparativamente as principais decisões metodológicas identificadas na literatura. Entre os estudos analisados, sete apresentam maior proximidade com as decisões investigadas nesta pesquisa: \citeonline{abasabadi2021automatic}, \citeonline{drotar2019ensemble}, \citeonline{fu2023reranking}, \citeonline{janani2023ensemble}, \citeonline{kaur2023mcdm}, \citeonline{seijopardo2019automatic} e \citeonline{zhang2019ensemble}. Esses trabalhos recebem maior aprofundamento ao longo da discussão, considerando aspectos relacionados à composição do \textit{ensemble}, à agregação dos resultados e à definição do subconjunto final de atributos.

%%%%%%%%%%%%%%%%%%%%%
%%%%%%%%%%%%%%%%%%%%%
\begin{sidewaystable}[p]
\centering
\caption{Síntese dos 22 estudos primários de \textit{Ensemble Feature Selection} analisados.}
\label{tab:efs}
\fontsize{8}{9.2}\selectfont
\setlength{\tabcolsep}{2.5pt}
\renewcommand{\arraystretch}{1.10}
\begin{tabular}{|P{2.3cm}|P{2.6cm}|P{2.3cm}|P{2.8cm}|P{3.0cm}|P{3.0cm}|P{2.6cm}|}
\hline
\textbf{Trabalho} &
\textbf{Composição do \textit{ensemble}} &
\textbf{Representação} &
\textbf{Agregação / consenso} &
\textbf{Diversidade / pesos / critérios adicionais} &
\textbf{Definição do subconjunto} &
\textbf{Avaliação} \\ \hline
\hline
\citeonline{abasabadi2021automatic} & Múltiplos ranqueadores & \textit{Rankings} & \textit{Fast non-dominated sorting} & Múltiplos ranqueadores (geração de diversidade) & Limiar automático & Desempenho e número de atributos \\ \hline
\citeonline{abellana2022ensemble} & Filtros heterogêneos & \textit{Scores}\slash \textit{rankings} & EDAS & Critérios heterogêneos & Retenção de 20\% dos atributos & AUC, redução dimensional e estabilidade \\ \hline
\citeonline{drotar2019ensemble} & Filtros heterogêneos & \textit{Rankings} & Plurality, STV, Borda e Borda ponderada & \textit{Mean-shift clustering} e similaridade entre ranqueadores & Top-$k$/número requerido de atributos & Relevância, estabilidade e classificação \\ \hline
\citeonline{fu2023reranking} & Filtros heterogêneos univariados e multivariados & \textit{Rankings} & TOPSIS e NSGA-III & Complementaridade uni\slash multivariada & Otimização multiobjetivo & Acurácia e F-measure \\ \hline
\citeonline{geng2026stability} & U-test, RFECV, LASSO e informação mútua & Subconjuntos\slash indicadores & Combinação linear ponderada & Pesos otimizados por Bayes e estabilidade entre \textit{folds} & Consenso de seleção e interseção entre \textit{folds} & Desempenho interno/externo e estabilidade \\ \hline
\citeonline{hashemi2021pareto} & Filtros heterogêneos & \textit{Rankings}\slash matriz de decisão & Pareto e ordenação não dominada & Relevância e redundância & Número de atributos definido experimentalmente & Acurácia, F-score e tempo \\ \hline
\citeonline{he2023ensemble} & Filtros heterogêneos & \textit{Rankings}\slash subconjuntos & Combinação ponderada & Pesos adaptativos por \textit{softmax} & Top-$k$ & Desempenho e adequação dos seletores \\ \hline
\citeonline{janani2023ensemble} & Filtros heterogêneos & \textit{Rankings} & IOWA, NWBM, BON-OWA e BON-IOWA & Ponderação por entropia & Top-$k$ & Acurácia, F-measure, precisão e revocação \\ \hline
\citeonline{joseph2025hybrid} & Filtros heterogêneos & \textit{Rankings} & Borda, MC4 e Random Dictator com distância de Kendall & Diversidade entre métodos de agregação & Etapa \textit{wrapper}/RFE & Desempenho preditivo \\ \hline
\citeonline{kaur2023mcdm} & $\chi^2$, correlação, InfoGain e ReliefF & \textit{Scores} & TOPSIS & Critérios heterogêneos & Seleção a partir do \textit{score} final & nMCC e G-measure \\ \hline
\citeonline{kavitha2025decision} & Filtros heterogêneos & \textit{Rankings}\slash matriz de preferência & TOPSIS, VIKOR, CODAS e variantes & Representação qROPHF & \textit{Ranking} seguido de seleção & Desempenho e significância estatística \\ \hline
\end{tabular}
\end{sidewaystable}

\begin{sidewaystable}[p]
\ContinuedFloat
\centering
\caption{Síntese dos 22 estudos primários de \textit{Ensemble Feature Selection} analisados (continuação).}
\fontsize{8}{9.2}\selectfont
\setlength{\tabcolsep}{2.5pt}
\renewcommand{\arraystretch}{1.10}
\begin{tabular}{|P{2.3cm}|P{2.6cm}|P{2.3cm}|P{2.8cm}|P{3.0cm}|P{3.0cm}|P{2.6cm}|}
\hline
\textbf{Trabalho} &
\textbf{Composição do \textit{ensemble}} &
\textbf{Representação} &
\textbf{Agregação / consenso} &
\textbf{Diversidade / pesos / critérios adicionais} &
\textbf{Definição do subconjunto} &
\textbf{Avaliação} \\ \hline
\hline
\citeonline{lin2024ensemble} & SU, FDR e ReliefF com complementaridade discriminante & Subconjuntos & Estratégia de agregação aprimorada & Complementaridade entre critérios & Subconjunto orientado pela capacidade classificatória & Desempenho e redução dimensional \\ \hline
\citeonline{loggenberg2026cross} & 18 métodos de FS baseados em PLS & \textit{Rankings}\slash frequência & \textit{Weighted Borda Count} & Pesos por SRD e desempenho no treinamento & Cotovelo adaptativo e refinamento top-15\% & Transferibilidade, desempenho e redução \\ \hline
\citeonline{ma2024ensemble} & Múltiplos critérios baseados em \textit{neighbourhood rough sets} & \textit{Scores} & Fusão multicritério & Critérios heterogêneos & Busca sequencial & Desempenho preditivo \\ \hline
\citeonline{mera2021framework} & Múltiplos métodos de FS & \textit{Rankings}\slash subconjuntos & Consenso por frequência & Configuração flexível & Limiar de consenso configurável & Desempenho e estabilidade \\ \hline
\citeonline{olabisi2025heterogeneous} & Filtros estatísticos heterogêneos & \textit{Rankings} & Borda & Complementaridade entre filtros & Limiar seguido de consolidação & Desempenho em detecção de \textit{phishing} \\ \hline
\citeonline{seijopardo2019automatic} & Ranqueadores em diferentes configurações & \textit{Rankings} & Desenhos CT e TC com diferentes combinações & Perturbação\slash configuração do \textit{ensemble} & Percentuais fixos ou limiar automático por complexidade & Erro de classificação e número de atributos \\ \hline
\citeonline{chinglensana2026ef} & mRMR, PCFS, JoMIC e MIFS-ND & \textit{Rankings}\slash subconjuntos & Consenso por frequência & Particionamento, refinamento iterativo e controle de redundância & Consenso entre ranqueadores e poda por correlação & Desempenho, redução e robustez \\ \hline
\citeonline{tang2023ensemble} & Filtros heterogêneos & \textit{Scores} & TOPSIS & Ponderação por entropia & \textit{Ranking} seguido de seleção & Desempenho preditivo \\ \hline
\citeonline{wu2023ecm} & Múltiplos seletores & Subconjuntos\slash votos & Coassociação e \textit{Feature Kernel} & Relações entre atributos & Núcleo inicial por coassociação; seleção até atingir $k$ atributos & Desempenho, robustez e estabilidade \\ \hline
\citeonline{zhang2019ensemble} & SAM, InfoGain, GeoDE e ReliefF sob \textit{bootstrap} & \textit{Rankings} & \textit{Rank Products} & Ponderação pela estabilidade & Percentual de atributos & AUC e estabilidade \\ \hline
\citeonline{zheng2019feature} & Tsallis e \textit{Markov blanket} sob múltiplas evidências & Subconjuntos\slash evidências & Dempster--Shafer & Ponderação por evidência & Resultado da fusão & Desempenho preditivo \\ \hline
\end{tabular}
\end{sidewaystable}
%%%%%%%%%%%%%%%%%%%%%
%%%%%%%%%%%%%%%%%%%%%

%%%%%%%%%%%%%%%%%%%%%
%%%%%%%%%%%%%%%%%%%%%
\subsection{Configuração do \textit{Ensemble} e Representação dos Resultados}
\label{subsec:efs-configuracao}

Uma das primeiras decisões na construção de um EFS refere-se à composição do \textit{ensemble} e à forma de introduzir diversidade entre os seletores. Em configurações heterogêneas, diferentes métodos de seleção de atributos são aplicados ao mesmo conjunto de treinamento, explorando critérios distintos de relevância, dependência ou redundância. Essa estratégia é adotada por \citeonline{drotar2019ensemble}, \citeonline{abellana2022ensemble}, \citeonline{kaur2023mcdm}, \citeonline{janani2023ensemble}, \citeonline{lin2024ensemble}, \citeonline{he2023ensemble} e \citeonline{joseph2025hybrid}, entre outros. A combinação de métodos com diferentes critérios busca explorar informações complementares e reduzir a dependência de um único seletor.

Entre os trabalhos que exploram essa configuração, \citeonline{drotar2019ensemble} combinam oito métodos de filtro e investigam a diversidade entre os resultados produzidos. Os autores propõem o agrupamento dos ranqueadores por meio do algoritmo \textit{mean-shift}, representando cada grupo por um \textit{ranking} médio antes da agregação. Essa estratégia busca reduzir a influência de seletores que apresentam resultados semelhantes, evitando que métodos redundantes tenham participação excessiva na formação do consenso. De forma relacionada, \citeonline{abasabadi2021automatic} propõem o \textit{Automatic Thresholding Feature Selection} (ATFS), que combina métodos de filtro e considera a similaridade entre os \textit{rankings} por meio de medidas de correlação, incluindo Spearman e Kendall. Embora empreguem mecanismos distintos, ambos os trabalhos destacam a importância de considerar a relação entre os seletores, e não apenas a quantidade de métodos combinados.

Outras configurações heterogêneas são investigadas por \citeonline{fu2023reranking} e \citeonline{janani2023ensemble}. Os primeiros combinam seis métodos de filtro, incluindo Chi-quadrado, Gini, Fisher Score, ReliefF, mRMR e CMIM, abrangendo diferentes critérios de avaliação dos atributos. Já \citeonline{janani2023ensemble} utilizam cinco métodos de seleção, incluindo abordagens baseadas em relevância e redundância, para investigar diferentes operadores de agregação. Em ambos os casos, a utilização de seletores heterogêneos permite reunir diferentes perspectivas sobre a importância dos atributos, embora o tratamento dessas informações na etapa de agregação seja distinto.

A diversidade também pode ser introduzida por meio de perturbações nos dados de treinamento. Essa estratégia é explorada por \citeonline{zhang2019ensemble}, que propõem o \textit{Ensemble Feature Selection Integrating Stability} (EFSIS). O método combina diferentes funções de seleção com amostras geradas por \textit{bootstrap}, permitindo avaliar tanto a relevância dos atributos quanto a consistência dos resultados diante de variações nos dados. Diferentemente das abordagens que exploram apenas a heterogeneidade dos seletores, o EFSIS incorpora a estabilidade como um componente adicional na construção do \textit{ensemble}.

Além da composição do \textit{ensemble}, a representação dos resultados produzidos pelos métodos de base influencia diretamente as possibilidades de combinação. Os seletores podem fornecer \textit{rankings}, \textit{scores} de relevância ou subconjuntos de atributos. A agregação de \textit{rankings} considera as posições relativas dos atributos, enquanto a combinação de \textit{scores} pode exigir procedimentos de normalização para tornar as avaliações comparáveis. Quando os resultados são subconjuntos, o consenso pode ser estabelecido a partir das decisões de inclusão ou exclusão dos atributos.

Nos estudos analisados, \citeonline{drotar2019ensemble}, \citeonline{joseph2025hybrid} e \citeonline{olabisi2025heterogeneous} utilizam predominantemente \textit{rankings}. Por outro lado, \citeonline{abellana2022ensemble}, \citeonline{kaur2023mcdm} e \citeonline{tang2023ensemble} organizam avaliações de relevância dos atributos como critérios de decisão multicritério. Em particular, \citeonline{kaur2023mcdm} estruturam os resultados de diferentes seletores em uma matriz de decisão, na qual os atributos representam as alternativas e as avaliações dos métodos correspondem aos critérios utilizados na agregação. Outras abordagens, como as de \citeonline{wu2023ecm}, \citeonline{mera2021framework} e \citeonline{zheng2019feature}, exploram o consenso entre subconjuntos ou evidências de seleção.

Essas diferentes configurações evidenciam que a composição do \textit{ensemble} e a representação de seus resultados são decisões relacionadas. A heterogeneidade dos métodos pode ampliar a complementaridade das avaliações, mas também introduzir redundância entre seletores. Da mesma forma, a utilização de posições, \textit{scores} ou decisões de seleção determina quais informações estarão disponíveis para a etapa de agregação e, consequentemente, quais mecanismos poderão ser empregados na obtenção do subconjunto final de atributos.

%%%%%%%%%%%%%%%%%%%%%
%%%%%%%%%%%%%%%%%%%%%
\subsection{Estratégias de Agregação, Ponderação e Consenso}
\label{subsec:efs-agregacao}

As estratégias de agregação identificadas nos estudos analisados abrangem métodos eleitorais, operadores de agregação, técnicas de decisão multicritério, abordagens de otimização e mecanismos de consenso estrutural. Essas estratégias diferem quanto à forma de combinar as avaliações produzidas pelos seletores, podendo utilizar posições em \textit{rankings}, \textit{scores} de relevância ou outras representações das informações sobre os atributos.

Entre os métodos eleitorais, \citeonline{drotar2019ensemble} investigam diferentes mecanismos de decisão coletiva, incluindo voto plural, voto único transferível, contagem de Borda e Borda ponderada. A proposta também considera o agrupamento prévio de ranqueadores, permitindo relacionar o tratamento da redundância entre seletores à formação do consenso. Outras aplicações de métodos eleitorais são apresentadas por \citeonline{olabisi2025heterogeneous}, que utilizam Borda para combinar filtros estatísticos na detecção de \textit{phishing}, e por \citeonline{joseph2025hybrid}, que investigam mecanismos como Borda, MC4 e distância de Kendall.

Além dos métodos eleitorais, diferentes operadores matemáticos têm sido empregados para combinar os resultados dos seletores. \citeonline{janani2023ensemble} investigam quatro operadores de agregação, incluindo IOWA e combinações baseadas em Bonferroni e OWA, utilizando os resultados de diferentes métodos de seleção de atributos. A proposta explora mecanismos de ponderação para produzir uma avaliação agregada, permitindo investigar o efeito da estratégia de combinação sobre os resultados do \textit{ensemble}. Embora os operadores utilizados sejam distintos de Borda, RRF e TOPSIS, o trabalho apresenta proximidade metodológica com o EFS4AC por tratar a escolha do mecanismo de agregação como uma decisão central.

Outra linha de investigação formula a agregação como um problema de decisão multicritério, no qual os atributos representam alternativas e as avaliações dos métodos de FS constituem os critérios de decisão. \citeonline{abellana2022ensemble} empregam EDAS, enquanto \citeonline{kaur2023mcdm} propõem o MCDM-EFS, que utiliza TOPSIS para combinar avaliações produzidas por $\chi^2$, correlação, \textit{Information Gain} e ReliefF. Nesse último trabalho, os \textit{scores} de importância são organizados em uma matriz de decisão, e o procedimento gera uma avaliação agregada utilizada para ordenar os atributos. \citeonline{tang2023ensemble} também adotam TOPSIS, incorporando ponderação baseada em entropia, enquanto \citeonline{kavitha2025decision} exploram TOPSIS, VIKOR e CODAS em um ambiente \textit{q-rung orthopair hesitant fuzzy}, buscando representar incertezas no processo de decisão.

A utilização de TOPSIS também é investigada por \citeonline{fu2023reranking}, que propõem o RTEFS, uma arquitetura de múltiplos estágios que combina informações provenientes de filtros univariados e multivariados. Além da agregação multicritério, a proposta incorpora procedimentos de reordenação dos atributos e otimização multiobjetivo por NSGA-III, considerando aspectos relacionados à relevância atributo-classe e à redundância atributo-atributo. Comparativamente ao MCDM-EFS, o RTEFS apresenta uma arquitetura mais complexa, na qual TOPSIS integra um processo mais amplo de refinamento das avaliações. Ambos os trabalhos são relevantes para o EFS4AC por investigarem a combinação de diferentes avaliações de atributos por meio de TOPSIS, embora não adotem necessariamente a mesma representação de entrada prevista nesta pesquisa.

Abordagens baseadas em otimização e ordenação multiobjetivo também são identificadas na literatura. \citeonline{hashemi2021pareto} utilizam dominância de Pareto e distância de aglomeração para ordenar atributos segundo critérios de relevância e redundância. De forma relacionada, \citeonline{abasabadi2021automatic} empregam \textit{fast non-dominated sorting} para integrar avaliações produzidas por métodos de filtro, associando esse procedimento à definição automática dos atributos selecionados. Diferentemente das abordagens baseadas em votação ou em uma função agregadora única, esses métodos consideram conjuntamente diferentes critérios de avaliação por meio de relações de dominância.

Outros estudos exploram mecanismos de consenso que não dependem exclusivamente da combinação de posições ou \textit{scores}. \citeonline{wu2023ecm} utilizam matrizes de coassociação para representar relações entre atributos; \citeonline{ma2024ensemble} combinam critérios derivados de \textit{neighbourhood rough sets}; e \citeonline{zheng2019feature} empregam a teoria de evidências de Dempster-Shafer para integrar subconjuntos. \citeonline{lin2024ensemble}, por sua vez, incorporam informações de complementaridade discriminante à combinação de diferentes filtros. O EF2RB, proposto por \citeonline{chinglensana2026ef}, utiliza mecanismos de consenso e refinamento iterativo, incluindo o controle da redundância entre os atributos selecionados.

Além da escolha do mecanismo de agregação, alguns estudos investigam a atribuição de pesos distintos aos componentes do \textit{ensemble}. \citeonline{he2023ensemble} utilizam uma função \textit{softmax} para ajustar adaptativamente os pesos dos seletores conforme sua adequação aos dados. \citeonline{zhang2019ensemble}, por sua vez, empregam \textit{Rank Products} para agregar os \textit{rankings} obtidos a partir de amostras geradas por \textit{bootstrap} e incorporam estimativas de estabilidade na ponderação dos diferentes métodos de seleção. Outras abordagens investigam pesos associados à confiabilidade ou à estabilidade dos componentes, como em \citeonline{loggenberg2026cross} e \citeonline{geng2026stability}.

Em conjunto, os estudos evidenciam que a agregação constitui uma decisão central na construção de um EFS, envolvendo a escolha do mecanismo de combinação e a possibilidade de ponderar diferentemente as contribuições dos seletores. Para o EFS4AC, essas alternativas são particularmente relevantes, pois permitem investigar como diferentes mecanismos de agregação influenciam o consenso produzido a partir dos \textit{rankings} individuais.

%%%%%%%%%%%%%%%%%%%%%
%%%%%%%%%%%%%%%%%%%%%
\subsection{Definição do Subconjunto Final e Estabilidade}
\label{subsec:efs-estabilidade-corte}

A definição do subconjunto final de atributos constitui uma etapa importante do EFS, podendo ocorrer após a agregação dos resultados dos seletores ou estar incorporada ao próprio mecanismo de consenso. Essa decisão influencia o resultado da seleção: um subconjunto excessivamente amplo pode manter atributos irrelevantes ou redundantes, enquanto uma seleção muito restritiva pode eliminar informações discriminativas.

Uma estratégia recorrente consiste em selecionar uma quantidade ou um percentual predefinido de atributos a partir dos resultados agregados. Essa abordagem é adotada por \citeonline{drotar2019ensemble}, \citeonline{janani2023ensemble}, \citeonline{he2023ensemble} e \citeonline{zhang2019ensemble}, entre outros. Em particular, \citeonline{drotar2019ensemble} selecionam uma quantidade requerida de atributos a partir dos resultados dos mecanismos eleitorais, enquanto \citeonline{janani2023ensemble} utilizam quantidades predefinidas para avaliar diferentes operadores de agregação. De forma semelhante, \citeonline{kaur2023mcdm} selecionam os atributos mais bem posicionados na ordenação produzida pelo TOPSIS. Embora essas estratégias sejam relativamente simples, a escolha da quantidade de atributos pode influenciar tanto o desempenho preditivo quanto o grau de redução da dimensionalidade.

Outros trabalhos investigam mecanismos que procuram determinar o subconjunto de maneira automática. \citeonline{seijopardo2019automatic} propõem a utilização de medidas de complexidade dos dados para estabelecer limiares de seleção em EFS, comparando essa estratégia com abordagens baseadas em quantidades ou percentuais predefinidos. Os autores também investigam duas configurações metodológicas: a combinação dos resultados dos seletores seguida da aplicação de limiares e a aplicação de limiares aos resultados individuais seguida da combinação. Essa distinção evidencia que a definição do subconjunto não depende apenas do valor de corte, mas também da ordem em que as operações de agregação e seleção são realizadas.

A determinação automática do subconjunto também é investigada por \citeonline{abasabadi2021automatic}, que propõem um procedimento baseado em métodos de filtro e \textit{fast non-dominated sorting}. Em vez de estabelecer previamente a quantidade de atributos, a abordagem procura identificar o subconjunto a partir das avaliações produzidas pelos seletores. Outras estratégias são apresentadas por \citeonline{mera2021framework}, que utilizam um limiar configurável sobre a frequência de seleção, e por \citeonline{wu2023ecm}, que iniciam a seleção a partir de um núcleo de atributos obtido por coassociação.

Em abordagens de múltiplos estágios, a definição do subconjunto pode estar associada a procedimentos adicionais de refinamento. \citeonline{fu2023reranking}, por exemplo, integram pré-seleção, reordenação e otimização na construção do subconjunto final. Essa organização difere das abordagens em que a agregação e a aplicação de uma regra de corte constituem etapas separadas, permitindo investigar individualmente seus efeitos sobre os atributos selecionados.

Além dos critérios de corte, a estabilidade constitui outra dimensão relevante na avaliação e construção de métodos de EFS. Essa propriedade expressa a consistência dos atributos selecionados diante de variações nos dados de treinamento. \citeonline{mera2021framework} avaliam a estabilidade utilizando o índice de Jaccard, enquanto \citeonline{drotar2019ensemble} analisam essa propriedade em conjunto com a identificação de atributos relevantes e o desempenho classificatório. No EFSIS, \citeonline{zhang2019ensemble} incorporam a estabilidade diretamente ao processo de agregação, atribuindo maior influência aos resultados considerados mais consistentes entre diferentes amostras. Assim, a estabilidade pode ser utilizada tanto como critério de avaliação quanto como componente do próprio método de seleção.

Em conjunto, os estudos mostram que a definição do subconjunto final pode envolver quantidades ou percentuais predefinidos, limiares automáticos e procedimentos de seleção integrados à agregação ou à otimização. Para o EFS4AC, essa distinção é relevante porque permite investigar separadamente os efeitos da agregação dos \textit{rankings} e dos critérios utilizados para selecionar os atributos que serão disponibilizados ao classificador associativo.


% =====================================================================
% =====================================================================
\section{Síntese das evidências e lacuna de pesquisa}
\label{sec:consideracoes-relacionados}
% =====================================================================

A análise desenvolvida neste capítulo permite relacionar duas vertentes da
literatura que, embora compartilhem o problema da seleção de atributos, apresentam
níveis distintos de desenvolvimento. A primeira compreende os estudos que aplicam
\textit{Feature Selection} (FS) diretamente à classificação associativa; a segunda
abrange as abordagens de \textit{Ensemble Feature Selection} (EFS), nas quais
múltiplas fontes de evidência são combinadas para reduzir a dependência de um
único método de seleção. A interseção entre essas duas vertentes delimita a lacuna
de pesquisa de interesse neste trabalho.

Na literatura de FS aplicada à classificação associativa, a seleção de atributos já
se encontra estabelecida como mecanismo para reduzir a dimensionalidade antes ou
durante a indução das regras. Conforme discutido na Seção~\ref{sec:fs-ca}, nove
dos 12 estudos analisados realizam a seleção de forma externa, isto é, produzem
primeiramente um subconjunto de atributos e somente depois executam o
classificador associativo. Os demais trabalhos apresentam formas de integração mais
estreitas, nas quais critérios de relevância participam da construção dos
antecedentes, da ponderação dos atributos ou da própria avaliação das regras.
Portanto, a lacuna identificada não corresponde à ausência de FS em classificação
associativa.

Os estudos analisados também não fornecem evidências para a adoção de um único
seletor como solução geral. Diferentes critérios informacionais, estatísticos,
baseados em modelos, em relevância e redundância ou em estratégias de busca
apresentam comportamentos distintos conforme o conjunto de dados e o
classificador considerado. Essa variabilidade é particularmente importante em
classificação associativa, pois a alteração do subconjunto de atributos modifica
não apenas a entrada do classificador, mas também o conjunto de itens disponível
para mineração e, consequentemente, o espaço no qual as Regras de Associação
Classificativas (RACs) podem ser geradas.

Nesse cenário, o trabalho de Sowan et al. (2023) constitui uma exceção relevante
dentro do conjunto analisado. Os autores utilizam o WARF, que combina CFS,
\textit{Information Gain}, \textit{Gain Ratio}, $\chi^2$ e ReliefF antes da
indução do PCBA, e propõem ainda uma extensão híbrida baseada em PSO. O estudo
demonstra, portanto, a viabilidade de associar um \textit{ensemble} de métodos de
FS a um classificador associativo. Entretanto, a proposta investiga uma
configuração específica de composição e agregação, orientada ao domínio de câncer
de mama, não tendo como objetivo comparar sistematicamente diferentes decisões de
projeto de um EFS, tais como mecanismo de consenso, ponderação dos seletores,
estabilidade da seleção e definição do ponto de corte.

Essa limitação torna-se mais evidente quando os resultados da Seção~\ref{ensambleFS}
são considerados. Os 22 estudos primários de EFS analisados apresentam diferentes
soluções para as principais etapas do processo. \citeonline{drotar2019ensemble},
por exemplo, investigam métodos eleitorais e diversidade entre ranqueadores;
\citeonline{zhang2019ensemble} incorporam estabilidade à ponderação dos
seletores; \citeonline{fu2023reranking} estruturam a agregação em múltiplos
estágios; \citeonline{loggenberg2026cross} associam \textit{Weighted Borda
Count}, ponderação e corte adaptativo; \citeonline{chinglensana2026ef} combinam
consenso, refinamento iterativo e controle de redundância; e
\citeonline{geng2026stability} relacionam otimização de pesos e
reprodutibilidade entre partições. Esses trabalhos demonstram que composição,
agregação, ponderação, estabilidade e corte são componentes independentes e
relevantes no projeto de um EFS.

Apesar desse avanço metodológico, nenhum dos seis estudos de maior proximidade
analisados na Seção~\ref{subsec:efs-sintese} investiga essas decisões no contexto
da classificação associativa. Suas avaliações são realizadas com classificadores
tradicionais e em domínios como bioinformática, categorização textual,
viticultura espectral e detecção de intrusões. Consequentemente, permanece em
aberto como estratégias desenvolvidas na literatura geral de EFS se comportam
quando o subconjunto selecionado é utilizado não apenas para treinar um
classificador convencional, mas para restringir um processo de mineração de
regras.

A lacuna identificada neste trabalho encontra-se, portanto, na interseção entre
essas duas vertentes. Entre os estudos analisados, permanece pouco explorada uma
investigação sistemática de EFS heterogêneo aplicado à classificação associativa
que permita analisar, sob um protocolo experimental comum, como diferentes
critérios de seleção e diferentes mecanismos de combinação afetam o subconjunto
de atributos utilizado na mineração. Particularmente, ainda requer investigação
o efeito conjunto de três decisões: a diversidade dos métodos de FS que compõem
o \textit{ensemble}, a estratégia utilizada para produzir o consenso e a regra
empregada para transformar esse consenso no subconjunto final de atributos.

Além disso, a avaliação dessa integração demanda critérios específicos da
classificação associativa. Em classificadores convencionais, um método de FS pode
ser avaliado principalmente pela redução dimensional e pelo desempenho
preditivo. Em classificação associativa, entretanto, a seleção repercute sobre
todo o processo de indução: reduz ou amplia o número de itens disponíveis,
modifica o espaço de busca de \textit{itemsets}, interfere na quantidade e na
estrutura das regras geradas e pode alterar o custo computacional e a
interpretabilidade do modelo. Dessa forma, uma avaliação adequada deve considerar,
além de medidas preditivas, propriedades como redução de atributos, estabilidade
dos subconjuntos, tempo de processamento, número de regras e complexidade das
RACs resultantes.

A estabilidade merece atenção particular nesse contexto. Uma estratégia de FS
pode preservar o desempenho médio do classificador e, ainda assim, selecionar
conjuntos substancialmente diferentes quando pequenas alterações são realizadas
nos dados de treinamento. Em classificação associativa, essa instabilidade pode
ser propagada para a etapa de mineração, produzindo conjuntos distintos de itens,
\textit{itemsets} e regras. Assim, avaliar a estabilidade do EFS não representa
apenas uma análise adicional da etapa de pré-processamento, mas também uma forma
de investigar a reprodutibilidade estrutural do classificador associativo obtido.

Dessa forma, a contribuição pretendida não parte da hipótese de que o EFS seja
inédito em classificação associativa, uma vez que existe precedente explícito na
literatura. O problema de pesquisa decorre da ausência, entre os trabalhos
analisados, de uma investigação abrangente das decisões que estruturam um EFS
nesse contexto e de seus efeitos sobre o processo associativo. Essa lacuna
fundamenta a continuidade desta pesquisa, voltada à análise de estratégias capazes
de combinar evidências produzidas por diferentes métodos de seleção antes da
indução do classificador associativo e avaliar seus efeitos tanto sobre a seleção
de atributos quanto sobre o desempenho e a estrutura do modelo baseado em regras.